\subsubsection{Convergence of Distributions}\label{subsec:distribution-convergence}

Just as sequences of classical functions are central to PDEs, so are sequences of generalized functions.

\begin{Definition}
\label{def:distribution-convergence}
A sequence $f_n$ in $\mathcal{D}^{\prime}(\Omega)$ \textbf{converges} to $f\in\mathcal{D}^{\prime}(\Omega)$ if \[(f_n,\phi)\to(f,\phi)\] for every $\phi\in\mathcal{D}(\Omega)$.
\end{Definition}

\begin{Example}
\label{ex:uniform-convergence-distribution}
A uniformly convergent sequence of continuous functions (which define distributions as in Example~\ref{ex:regular-distribution}) also converges in $\mathcal{D}^{\prime}$.
\end{Example}

\begin{Example}
\label{ex:delta-rectangular-approximation}
Consider the sequence
\begin{equation}\label{eq:delta-approximating-rectangular-sequence}
f_n(x)=\begin{cases}n,&0<x<1/n\\0,&\text{otherwise.}\end{cases}
\end{equation}
We have
\begin{equation}\label{eq:delta-rectangular-sequence-pairing}
\int_{-\infty}^{\infty}f_n(x)\phi(x)\,\mathrm{d}x=n\int_0^{1/n}\phi(x)\,\mathrm{d}x,
\end{equation}
which converges to $\phi(0)$ as $n\to\infty$. Hence $f_n(x)\to\delta(x)$ in $\mathcal{D}^{\prime}(\mathbb{R})$.
\end{Example}

\begin{Remark}
\label{rem:distribution-density}
Problem~\ref{ex:compactly-supported-distribution-density} asks the reader to prove that every distribution is the limit of distributions with compact support. Later we shall actually see that every distribution is a limit of test functions; in other words, test functions are dense in $\mathcal{D}^{\prime}(\Omega)$.
\end{Remark}

Another important result is the (sequential) completeness of $\mathcal{D}^{\prime}(\Omega)$.

\begin{Theorem}
\label{thm:distribution-sequential-completeness}
Let $f_n$ be a sequence in $\mathcal{D}^{\prime}(\Omega)$ such that $(f_n,\phi)$ converges for every $\phi\in\mathcal{D}(\Omega)$. Then there exists $f\in\mathcal{D}^{\prime}(\Omega)$ such that $f_n\to f$.
\end{Theorem}

\begin{Proof}
We define
\begin{equation}\label{eq:distribution-pointwise-pairing-limit}
(f,\phi)=\lim_{n\to\infty}(f_n,\phi).
\end{equation}
Obviously, $f$ is a linear mapping from $\mathcal{D}(\Omega)$ to $\mathbb{R}$. To verify that $f\in\mathcal{D}^{\prime}(\Omega)$, we have to establish its continuity, i.e., we must show that if $\phi_n\to0$ in $\mathcal{D}(\Omega)$, then $(f,\phi_n)\to0$. Assume the contrary. Then, after choosing a subsequence which we again label $\phi_n$,\footnote{The use of the same symbol for both a sequence and any of its subsequences is a typical practice in PDEs. Its primary purpose is clarity of notation (since we often have to consider subsequences several levels deep), but it has the pleasant side effect of driving many classical analysts crazy. Of course, there are cases where it is important to distinguish between a sequence and its subsequence (as we do later in this proof) and we do so with appropriate notation.} we may assume $\phi_n\to0$, but $|(f,\phi_n)|\geq c>0$.

Now recall that convergence to 0 in $\mathcal{D}(\Omega)$ means that the supports of all the $\phi_n$ lie in a fixed compact subset of $\Omega$ and that all derivatives of the $\phi_n$ converge to zero uniformly. After again choosing a subsequence, we may assume that $|D^\alpha\phi_n(\bm{x})|\leq4^{-n}$ for $|\alpha|\leq n$. Let now $\psi_n:=2^n\phi_n$. Then the $\psi_n$ still converge to 0 in $\mathcal{D}(\Omega)$, but $|(f,\psi_n)|\to\infty$.

We shall now recursively construct a subsequence $\{f_n'\}$ of $\{f_n\}$ and a subsequence $\{\psi_n'\}$ of $\{\psi_n\}$. First we choose $\psi_1'$ such that $|(f,\psi_1')|>1$. Since $(f_n,\psi_1')\to(f,\psi_1')$, we may choose $f_1'$ such that $|(f_1',\psi_1')|>1$. Now suppose we have chosen $f_j'$ and $\psi_j'$ for $j<n$. We then choose $\psi_n'$ from the sequence $\{\psi_n\}$ such that
\begin{equation}\label{eq:distribution-completeness-small-earlier-pairings}
|(f_j',\psi_n')|<\dfrac1{2^{n-j}},\quad j=1,2,\ldots,n-1,
\end{equation}

\begin{equation}\label{eq:distribution-completeness-large-limit-pairing}
|(f,\psi_n')|>\sum_{j=1}^{n-1}|(f,\psi_j')|+n.
\end{equation}
This is possible because, on the one hand, $\psi_n\to0$, and, on the other hand, $|(f,\psi_n)|\to\infty$. Since, moreover, $(f_n,\psi)\to(f,\psi)$, we can choose $f_n'$ such that
\begin{equation}\label{eq:distribution-completeness-diagonal-pairing}
|(f_n',\psi_n')|>\sum_{j=1}^{n-1}|(f_n',\psi_j')|+n.
\end{equation}
Next we set
\begin{equation}\label{eq:distribution-completeness-test-series}
\psi=\sum_{n=1}^{\infty}\psi_n'.
\end{equation}
It follows from the construction of the $\psi_n'$ that the series on the right converges in $\mathcal{D}(\Omega)$. Hence
\begin{equation}\label{eq:distribution-completeness-pairing-series-split}
(f_n',\psi)=\sum_{j=1}^{n-1}(f_n',\psi_j')+(f_n',\psi_n')+\sum_{j=n+1}^{\infty}(f_n',\psi_j').
\end{equation}
From \eqref{eq:distribution-completeness-small-earlier-pairings} we find that
\begin{equation}\label{eq:distribution-completeness-tail-estimate}
\left|\sum_{j=n+1}^{\infty}(f_n',\psi_j')\right|<\sum_{j=n+1}^{\infty}2^{n-j}=1,
\end{equation}
and this in conjunction with \eqref{eq:distribution-completeness-pairing-series-split} and \eqref{eq:distribution-completeness-diagonal-pairing} implies that $|(f_n',\psi)|>n-1$. Hence the limit of $(f_n',\psi)$ as $n\to\infty$ does not exist, a contradiction.
\end{Proof}

A similar contradiction argument can be used to prove the following lemma; the details of the proof are left as an exercise (cf. Problem~\ref{ex:distribution-test-sequence-pairing}).

\begin{Lemma}
\label{lem:distribution-test-sequence-pairing}
Assume that $f_n\to0$ in $\mathcal{D}^{\prime}(\Omega)$ and $\phi_n\to0$ in $\mathcal{D}(\Omega)$. Then $(f_n,\phi_n)\to0$.
\end{Lemma}

We also have the following corollary.

\begin{Corollary}
\label{cor:distribution-joint-continuity}
If $f_n\to f$ and $\phi_n\to\phi$, then $(f_n,\phi_n)\to(f,\phi)$.
\end{Corollary}

Hence the pairing between distributions and test functions is continuous. (Of course separate continuity in each factor is obvious from the definitions, but joint continuity requires a proof.)

\begin{Proof}
The corollary follows immediately from the identity
\begin{equation}\label{eq:distribution-joint-pairing-identity}
(f_n-f,\phi_n-\phi)=(f_n,\phi_n)-(f,\phi_n)-(f_n,\phi)+(f,\phi).
\end{equation}

\end{Proof}
